\clearpage
\begin{appendices}
\renewcommand{\thesubsection}{\Alph{section}.\arabic{subsection}}
\section{支撑材料与运行方法}
\label{sec:appendix-materials}
\subsection{文件清单与来源}
表\ref{tab:files}列出论文支撑材料。题目及附件1、附件2、附件3模板为官方输入，按原目录保存。附录B完整列出正式计算、数值核验、作图和数据导出的10个Python源文件，库函数本体不重复附入。当前稿图件及绘图源文件随支撑材料保存。

\begin{longtable}{@{}>{\raggedright\arraybackslash}p{0.38\textwidth}>{\raggedright\arraybackslash}p{0.57\textwidth}@{}}
\caption{支撑材料文件清单}\label{tab:files}\\
\toprule 文件或相对路径&作用\\\midrule\endfirsthead
\toprule 文件或相对路径&作用\\\midrule\endhead
\bottomrule\endfoot
\texttt{code/common.py}&环境插值、圆柱网格、有限体积组装、追赶法及时间推进\\
\texttt{code/p1.py}&问题一计算、五级细化和结果输出\\
\texttt{code/p23.py}&问题二、三耦合求解、收敛分析与事件定位\\
\texttt{code/p4.py}&收缩域求解、结构对照、物理位置重采样\\
\texttt{code/verify.py}&半解析温度、非线性制造解与线性求解器交叉核验\\
\texttt{code/sensitivity.py}&两类长期模型的14组参数扰动计算\\
\texttt{code/figures.py}&九幅计算结果与检验图件的绘制，含当前图\ref{fig:p34-comparison}的字体、图例及标注设置\\
\texttt{code/audit.py}&工作簿与原稿六张结果表核对、四问最终计算设置复算\\
\texttt{code/}\allowbreak\texttt{export\_figure2\_data.py}&从附件1导出图\ref{fig:p1-environment}的241条环境记录，并计算绘图坐标范围\\
\texttt{code/}\allowbreak\texttt{export\_figure8b\_data.py}&从问题四未舍入结果和附件2导出图\ref{fig:p4-drying}(b)的绘图数据\\
结果/\texttt{result1.xlsx}--\texttt{result4.xlsx}&四问按指定模板与四位小数格式保存的结果\\
结果/\texttt{p1\_field.npz}--\texttt{p4\_field.npz}&未舍入采样场、问题三及四的临界时刻等数据\\
结果/\texttt{sensitivity.npz}&参数扰动得到的干燥时长\\
结果/\texttt{p1\_论文表格.md}等四个文件&六张指定表的原始数值文本\\
结果/\texttt{p1\_检验数据.md}等四个文件&既有网格、时间步和方程收支检验记录\\
结果/灵敏度分析.md&参数扰动数值记录，正文对部分早期解释作了限定\\
结果/\texttt{audit\_report.json}&一致性检查及计算复核记录\\
\texttt{paper/figures/}下的\texttt{图1.tex}&当前图\ref{fig:framework}的TikZ绘图源文件\\
\texttt{paper/figures/}下的\texttt{图2.png}&图\ref{fig:p1-schematic}的问题一模型示意图\\
\texttt{paper/figures/}下的\texttt{图9.png}&图\ref{fig:p4-schematic}的问题四收缩区域与固定物质坐标映射示意图\\
\texttt{paper/figures/}下的\texttt{图3.pdf}、\texttt{图11.pdf}&当前图\ref{fig:p1-environment}和图\ref{fig:p34-comparison}的矢量图件\\
\texttt{paper/figures/}下的\texttt{fig01\_env\_style.tex}、\texttt{fig01\_env\_records.dat}、\texttt{fig01\_env\_limits.tex}&图\ref{fig:p1-environment}的PGFPlots绘图源文件、逐分钟环境记录及坐标范围；源文件设置字体、坐标标签、图例和恒定延拓\\
\texttt{paper/figures/}下的图4(a).png、图4(b).png等10个文件&当前图\ref{fig:p1-temperature}、图\ref{fig:p1-moisture}、图\ref{fig:p2-evolution}、图\ref{fig:p3-drying}、图\ref{fig:p4-drying}的五组子图\\
\texttt{paper/figures/}下的\texttt{图12.png}&当前图\ref{fig:parameter-sensitivity}的灵敏度分析图\\
\texttt{paper/figures/}下的\texttt{fig08b\_style.tex}、\texttt{fig08b\_curves.dat}、\texttt{fig08b\_measurements.dat}、\texttt{fig08b\_endpoint.tex}&图\ref{fig:p4-drying}(b)的PGFPlots绘图源文件、曲线数据、半径记录采样点及临界时刻\\
\texttt{paper/67A论文.tex}、\texttt{论文正文.tex}、\texttt{论文附录.tex}及\texttt{sec/}目录&当前论文入口、正文入口、附录和分节内容；图\ref{fig:p2-coupling}的TikZ源码位于“问题二模型建立.tex”内\\
\texttt{paper/ai-usage.tex}&AI使用详情的可编辑源文件\\
\texttt{paper/AI工具使用详情.pdf}&单独提交的AI工具、过程及核验状态说明\\
\end{longtable}

\subsection{运行环境与复现顺序}
使用Python 3.12及NumPy、pandas、SciPy、openpyxl、Matplotlib，依赖版本见\texttt{requirements.txt}。后两者分别承担工作簿读写及作图，SciPy只提供特殊函数、标量求根和带状线性系统求解等基础数值功能。自编数值组装与控制流程参见后续源码。浮点计算用双精度；工作簿仅在最终写入时舍入，不将四位小数值反馈给求解器。

将官方“附件”文件夹放在项目根目录，在该目录依次执行下列命令。其中\texttt{p23.py p3}明确运行第三问；不带参数时运行第二问。主程序会重新生成“结果”目录内对应文件，原始附件模板不会被覆盖。\texttt{audit.py --solve}只重新计算并比较既有结果，不覆盖工作簿和采样场。
\begin{lstlisting}[language=bash,basicstyle=\ttfamily\small]
python -m pip install -r requirements.txt
python code/p1.py
python code/verify.py
python code/p23.py
python code/p23.py p3
python code/p4.py
python code/sensitivity.py
python code/figures.py
python code/audit.py --solve
\end{lstlisting}
\texttt{figures.py}生成计算图件，其中\texttt{fig\_compare()}按当前图\ref{fig:p34-comparison}的样式绘制两问的中心含水率曲线。当前稿使用表\ref{tab:files}中保存的PDF、PNG和TikZ源文件，其中图\ref{fig:framework}与图\ref{fig:p2-coupling}随论文编译；直接编译当前稿无需重新绘图。图\ref{fig:p1-environment}与图\ref{fig:p4-drying}(b)的重绘使用PGFPlots及XeLaTeX，中文和西文字体分别为SimSun和Times New Roman。

重绘图\ref{fig:p1-environment}时，先在项目根目录运行\texttt{python code/export\_figure2\_data.py}，再进入\texttt{paper/figures}目录执行以下命令，直接生成论文引用的\texttt{图3.pdf}。若已运行\texttt{figures.py}，应随后执行本步骤，以保留当前图\ref{fig:p1-environment}的字体、坐标标签和图例设置。
\noindent\begin{minipage}{\linewidth}
\begin{lstlisting}[language=bash,basicstyle=\ttfamily\small]
xelatex -interaction=nonstopmode -halt-on-error -jobname=图3 fig01_env_style.tex
\end{lstlisting}
\end{minipage}
重绘图\ref{fig:p4-drying}(b)时，先在项目根目录运行\texttt{python code/export\_figure8b\_data.py}，再在\texttt{paper/figures}目录用XeLaTeX编译\texttt{fig08b\_style.tex}，将生成的PDF按原始比例转为图10(b).png，图像尺寸为$2391\times2298$像素。

在\texttt{paper}目录执行下列编译命令两次，以解析正文与附录的交叉引用：
\begin{lstlisting}[language=bash,basicstyle=\ttfamily\small]
xelatex -interaction=nonstopmode -halt-on-error 67A论文.tex
\end{lstlisting}
另编译\texttt{ai-usage.tex}，将输出命名为“AI工具使用详情.pdf”。所有附录源码以\texttt{lstinputlisting}直接读入实际运行文件，使论文附录与支撑程序保持一致。

\subsection{输出与达标判定约定}
\texttt{result1.xlsx}和\texttt{result2.xlsx}分别含1800、10800个逐秒记录，各设“温度”“水分浓度”两张工作表。\texttt{result3.xlsx}和\texttt{result4.xlsx}分别含3449、3066个逐分钟记录，时间范围为60--206940 s和60--183960 s，均延伸至临界事件之后的首个整分钟。四个工作簿的固定径向采样位置均为0--2 cm、间隔0.1 cm；问题四另设末列“药材表面”，超出当前半径的单元留空。

所有场量以双精度计算，在写入工作簿时统一保留四位小数。严格达标条件依据未舍入含水率判断，末行显示0.1500仍可能对应小于0.15的计算值。问题三、四按式\eqref{eq:p3-event-interpolation}在首次跨过阈值的时间步内线性插值，正文表格的结束行采用临界时刻插值状态，与工作簿末行的整分钟状态分别对应不同时间。

\texttt{audit.py}逐单元核对工作簿舍入值、时间列、距离表头和域外空白；其六张结果表核验读取原稿\texttt{paper/body.tex}，当前稿沿用这些结果数值，并已另行与原始数值文本逐表核对。\texttt{--solve}选项进一步在相同设置下重算四问，比较完整采样场及临界时刻。

\subsection{加密网格的构造}
问题一至三采用外侧$b=0.5$ mm加密、内区均匀划分的径向网格。问题四将该网格按初始半径归一化，固定物质坐标中的加密带宽为0.025，对应的物理带宽随$R(t)$收缩。以下规则给出物理网格（问题四为初态网格）加密带的可复现设置，各细化等级由同一规则生成。
最细网格的表面最小步长名义值取$1.25\times10^{-7}$ m，几何比为$1.3^{1/16}$，再按加密带总宽度归一化。设带宽为$b$、名义最小步长为$h_{\min}$，加密区分段数及步长分别为
\begin{equation}
\label{eq:appendix-grid}
m=\left\lceil\frac{\log[1+b(q-1)/h_{\min}]}{\log q}\right\rceil,
\qquad \Delta r_j=\frac{bq^{m-j}}{\sum_{\ell=1}^m q^{m-\ell}},\quad j=1,\ldots,m.
\end{equation}

\clearpage
\section{完整程序源码}
\label{sec:appendix-code}
\begingroup
\lstset{language=Python,basicstyle=\ttfamily\fontsize{8}{10}\selectfont,
  commentstyle=\color{black!65},stringstyle=\color{black!80},
  keywordstyle=\bfseries,backgroundcolor=\color{white},
  frame=single,numbersep=5pt,xleftmargin=1.8em,framexleftmargin=1.5em,
  aboveskip=0.5em,belowskip=0.8em,breaklines=true,columns=fullflexible}
\lstset{literate={∈}{{$\in$}}1 {℃}{{$^{\circ}\mathrm{C}$}}1}
\subsection{公共数值模块：common.py}
\lstinputlisting{../code/common.py}
\subsection{问题一：p1.py}
\lstinputlisting{../code/p1.py}
\subsection{问题二与问题三：p23.py}
\lstinputlisting{../code/p23.py}
\subsection{问题四：p4.py}
\lstinputlisting{../code/p4.py}
\subsection{独立算法核验：verify.py}
\lstinputlisting{../code/verify.py}
\subsection{参数灵敏度：sensitivity.py}
\lstinputlisting{../code/sensitivity.py}
\subsection{图件生成：figures.py}
\lstinputlisting{../code/figures.py}
\subsection{结果一致性与计算复核：audit.py}
\lstinputlisting{../code/audit.py}
\subsection{图\ref{fig:p1-environment}数据导出：export\_figure2\_data.py}
\lstinputlisting{../code/export_figure2_data.py}
\subsection{图\ref{fig:p4-drying}(b)数据导出：export\_figure8b\_data.py}
\lstinputlisting{../code/export_figure8b_data.py}
\endgroup
\end{appendices}
